\section*{Senior Applied Scientist (AS) | Amazon, Weblab}
\textbf{September 2021 - Present} (\emph{09/21--09/22: AS I; 09/22--05/24: AS II; 05/24--current: AS III})

I am a Senior Applied Scientist at Weblab, Amazon’s largest A/B testing platform, where I work on advancing online experimentation at scale. My work spans across open-ended research, academic paper writing, and the hands-on implementation of large-scale algorithms in production, directly impacting hundreds of millions of customers daily. 

\par My main focus in Weblab has been:

\begin{itemize}
    \item \textbf{Online Experimentation with Interference:} 
    I have pioneered experimental designs for large-scale testing of interference in online experiments, co-leading the implementation of the first large ``multi-randomized'' experiments at Amazon. These experiments, unlike classical A/B tests, jointly randomize across multiple units in the experiment (e.g., customers and advertisers). Via these designs, we can detect ``spillover effects'' that would otherwise be overlooked. For example, these have been implemented to precisely capture the indirect behavioral response of advertisers induced by a change in pricing schemes of ad-auctions. Key papers I have authored include (asterisk denotes first authorship, double asterisk first co-authorship): 
    \small{
        \begin{itemize}
	    \item \href{https://doi.org/10.1093/jrsssb/qkaf073}{\color{blue}{Multiple Randomization Designs*}} (JRSS-B, 2026)
            \item {\href{https://projecteuclid.org/journals/statistical-science/advance-publication/Experimental-Design-in-Marketplaces/10.1214/23-STS883.short}{\color{blue}{Experimental Design in Marketplaces*}}} (Statistical Science, 2023)
            \item \href{https://www.amazon.science/publications/efficient-switchback-experiments-via-multiple-randomization-designs}{\color{blue}{Efficient Switchback Experiments via Multiple Randomization Designs*}} (CODE@MIT 2023)
            \item \href{https://arxiv.org/abs/2405.18621}{\color{blue}{Multi-Armed Bandits with Network Interference}} (NeurIPS 2024, and CODE@MIT 2024)
            \item \href{}{\color{blue}{Regression Adjustments for Experimental Designs in Two-Side Marketplaces}} (CODE@MIT 2024)
            \item \href{}{\color{blue}{Measuring direct and Indirect Impacts in a Multi-Sided Marketplace: Evidence from a Clustered Multiple Randomization Experiment**}} (TSMO@KDD 2024, CODE@MIT 2024)
        \end{itemize}
	}
	\normalsize{    
    \item \textbf{Enhancing Causal Estimates via Covariate Adjustments and Traffic Prediction:}
    A/B tests are notoriously noisy. To better estimate the impact of interventions, we can leverage unit-level side information (covariates) and embed it within regression models. 
    Over the last year, I lead the design and implementation of a distributed analytic engine for the analysis of very large volumes of experimental data produced by online experiments.
    Our system leverages a distributed architecture that heavily relies on PySpark for data processing. 
    Our product enables more accurate causal estimates in thousands of daily experiments. 
     This innovation was backed by research on the value of covariates in online A/B testing, contributing to improvements in experiment analytics.
     As part of this work stream, I also developed methods to better predict future traffic in online experiments. Featured papers (asterisk denotes first authorship, double asterisk first co-authorship):}
     \small{
    \begin{itemize}
        \item  \href{https://proceedings.mlr.press/v218/masoero23a.html}{\color{blue}{Leveraging Covariate Adjustments at Scale in Online A/B Testing*}} (KDD CDPD, 2023).
        \item \href{https://arxiv.org/pdf/2402.03231}{\color{blue}{Improved prediction of future online activity in online A/B testing*}} (in submission)
        \item \href{https://arxiv.org/abs/2401.14722}{\color{blue}{A Nonparametric Bayes Approach to Online Activity Prediction*}} (in submission)
    \end{itemize}
    }
\end{itemize}

Previously, I completed my PhD at MIT in EECS in August 2021 under the supervision of Professor Tamara Broderick. 
During my PhD, I worked on Bayesian inference, with an emphasis on the development of efficient and theoretically grounded approaches for optimal design of experiments in genomic studies. In particular, my PhD thesis consisted of a sequence of papers developing a methodology to devise optimal sampling strategies to maximize learnings stemming from rare-variants discoveries under fixed experimental budgets  using a Bayesian nonparametric approach (``\href{https://academic.oup.com/biomet/article-abstract/109/1/17/6146908?redirectedFrom=fulltext}{\color{blue}{More for Less: Predicting and maximizing genetic variant discovery via Bayesian nonparametrics}}'', in \href{https://academic.oup.com/biomet/advance-article-abstract/doi/10.1093/biomet/asab012/6146908?redirectedFrom=PDF}{Biometrika}, ``\href{https://www.tandfonline.com/doi/full/10.1080/01621459.2022.2115918}{\color{blue}{Scaled process priors: Improved predictions and uncertainties for new-feature counts via random scaling in Bayesian nonparametrics}}'', in the \href{https://www.tandfonline.com/doi/full/10.1080/01621459.2022.2115918}{Journal of the American Statistical Association},  ``\href{https://arxiv.org/abs/2112.02032}{\color{blue}{Bayesian nonparametric strategies for power maximization in rare variants association studies}}'', spotilight at \href{https://www.lmrl.org/}{Learning Meaningful Representations of Life},  Neural Information Processing Systems, 2021).